problem:  
Let \( M^n \) be a closed, simply connected Riemannian manifold of nonnegative sectional curvature, and suppose that \( T^k \) acts effectively and isometrically on \( M \) with \( k = \mathrm{srk}(M) \) (the symmetry rank of \( M \)) maximal among manifolds of nonnegative curvature in dimension \( n\). Known classification results (Grove–Searle, Wilking, Grove–Wilking–Ziller) imply that in the positively curved case, such \( M \) is diffeomorphic to a compact rank-one symmetric space (CROSS) or a biquotient thereof.  

The question is:  
> **Conjectural classification problem:**  
> If \( M \) has nonnegative sectional curvature and maximal symmetry rank, must \( M \) be diffeomorphic to a biquotient of a product of compact rank-one symmetric spaces (e.g. \( S^n \), \( \mathbb{CP}^n \), \( \mathbb{HP}^n \), or \( \mathrm{CaP}^2 \)) by a free torus action?  
> Equivalently, does the combination of nonnegative curvature and maximal torus symmetry determine \( M \) up to diffeomorphism within the class of \((X_1 \times \cdots \times X_r)/T^s\), where each \( X_i \) is a CROSS?

Why is it a "good" problem:  
This problem directly targets one of the central frontiers in Riemannian geometry and topology—the **classification of nonnegatively curved manifolds with large symmetry**. It proposes a sharp extension of the Grove–Searle and Wilking results from the positively curved case to the more flexible but far less understood nonnegative curvature regime.

It is a “good” problem because:  

1. **Conceptual depth and foresight:**  
   It asks whether *symmetry and curvature alone* rigidly constrain the topology of manifolds beyond known cases. A positive answer would unify several decades of geometric classification theorems under one schema and complete the analogy between the positive and nonnegative settings.  

2. **Bridges multiple fields:**  
   The problem links **global Riemannian geometry** (sectional curvature), **transformation group theory** (isometric torus actions and symmetry rank), and **differential topology** (classification through biquotient structures). Its resolution likely requires new insights combining Alexandrov geometry, rational homotopy theory, and the structure of torus manifolds.

3. **Mathematical simplicity masking great depth:**  
   Expressed in a single line—“Are all nonnegatively curved manifolds with maximal torus symmetry biquotients of CROSS products?”—the problem is deceptively simple, yet encodes the complexity of manifold topology under curvature constraints. It connects with active programs in topology (torus manifolds, rational ellipticity), geometry (soul theorem, orbit space structures), and algebraic groups (biquotient constructions).

4. **Positioned beyond existing results but plausible:**  
   Known classification theorems cover the positively curved case and many low-dimensional or cohomogeneity-one examples; however, no general result is known in nonnegative curvature. The conjecture’s formulation refines current understanding without contradicting any established theorem.

In short, this problem identifies a precise, open, and interfield conjectural frontier: the geometric–topological classification of **highly symmetric, nonnegatively curved manifolds**, a cornerstone question whose resolution would fundamentally advance global Riemannian geometry.